% CS 418: Intro to Data Science — Worksheet 9 (09-week05-day1)
% Week 5, Monday 2026-09-21
% Compile with: latexmk -pdf worksheet.tex

\documentclass[11pt]{article}

% ── Packages ──────────────────────────────────────────────────────────────────
\usepackage[margin=0.65in, headheight=14pt]{geometry}
\usepackage{amsmath, amssymb}
\usepackage{ragged2e}
\usepackage{enumitem}
\usepackage{booktabs}
\usepackage{graphicx}
\usepackage{hyperref}
\usepackage{xcolor}
\usepackage{fancyhdr}
\usepackage{mdframed}
\usepackage{array}


% ── Page style ────────────────────────────────────────────────────────────────
\pagestyle{fancy}
\fancyhf{}
\lhead{\textbf{CS 418: Intro to Data Science}}
\rhead{UIC, Fall 2026}
\cfoot{\thepage}
\renewcommand{\headrulewidth}{1pt}
\setlength{\headheight}{12pt}
\setlength{\headsep}{0.4cm}

% ── Hyperlinks ────────────────────────────────────────────────────────────────
\hypersetup{colorlinks=false, hidelinks}

% ── Convenience environments ──────────────────────────────────────────────────
% Usage: \blankspace{6cm} — blank area for hand-written work, no rules
\newcommand{\blankspace}[1]{%
  \vspace{0.3em}
  \begin{minipage}[t][#1][t]{\linewidth}
  \end{minipage}
  \vspace{0.3em}
}

% Usage: \guidedopen{title}{guided content}{guided workspace height}{open-ended workspace height}
% Left: a titled, structured prompt with its own blank workspace. Right: a
% fully blank workspace for open-ended exploration of the same topic (no
% prompt text), sized so the two columns land at roughly the same total
% height. Neither workspace is ruled.
\newcommand{\guidedopen}[4]{%
  \noindent\textbf{#1}\\[0.18em]
  \noindent
  \begin{minipage}[t]{0.485\textwidth}
    #2
    \blankspace{#3}
  \end{minipage}
  \hfill
  \begin{minipage}[t]{0.485\textwidth}
    \blankspace{#4}
  \end{minipage}
  \\[0.4em]
}

% Usage: \panelbox{width}{height}{caption} — a framed drawing panel with `caption`
% centered beneath it. An empty caption leaves the panel bare, with nothing
% below the frame at all.
\newcommand{\panelbox}[3]{%
  \begin{minipage}[t]{\dimexpr#1+2\fboxsep+2\fboxrule\relax}
    \framebox[#1][c]{\rule{0pt}{#2}}%
    \if\relax\detokenize{#3}\relax\else
      \\[0.25em]{\centering\small #3\par}%
    \fi
  \end{minipage}%
}

% ══════════════════════════════════════════════════════════════════════════════
\begin{document}
\RaggedRight

% ── Title block: topic left, info box right ───────────────────────────────────
\noindent
\begin{minipage}[t]{0.55\textwidth}
  \vspace{0pt}
  {\LARGE \textbf{Worksheet \#\,9}} \\[0.18em]
  {\large \textit{Week 5, Day 1: Elements of visualization, choosing a plot, correlation, perception}}
\end{minipage}
\hfill
\begin{minipage}[t]{0.40\textwidth}
  \vspace{0pt}
  \small
  Name: \underline{\hspace{4.5cm}} \\[0.35em]
  NetID: \underline{\hspace{4.3cm}} \\[0.35em]
  Date: \underline{\hspace{4.4cm}} \\
\end{minipage}

\vspace{0.6em}
\noindent\rule{\linewidth}{0.6pt}
\vspace{0.1em}

\noindent\underline{\textbf{Guided checkpoints (optional)}}
\vspace{0.8em}

% ── 1. Bertin's seven ────────────────────────────────────────────────────────
\guidedopen{1. Bertin's Seven Visual Variables}{
List all seven variables.
}{4.4cm}{5.3cm}

\vfill

% ── 2. Munzner's channels ────────────────────────────────────────────────────
\guidedopen{2. Munzner's Magnitude Channels}{
List the ``top five'' magnitude channels
}{4.4cm}{5.3cm}

\vfill

% ── 3. Choosing a plot ───────────────────────────────────────────────────────
\guidedopen{3. Choosing a Plot: One Quantitative Feature}{
Name three plots for a single quantitative feature. 
}{4.3cm}{5.4cm}

\vfill

\newpage

% ── 4. Correlation ───────────────────────────────────────────────────────────
\guidedopen{4. Pearson's Correlation Coefficient}{
Sketch a scatter plot for each panel.

\vspace{0.45em}

\noindent
\panelbox{2.3cm}{2.0cm}{$r \approx 0.9$}\hfill
\panelbox{2.3cm}{2.0cm}{$r \approx 0$}\hfill
\panelbox{2.3cm}{2.0cm}{$r \approx -0.5$}
}{0.4cm}{3.9cm}

\vfill

% ── 5. Simpson's paradox ─────────────────────────────────────────────────────
\guidedopen{5. Simpson's Paradox}{
Sketch the same data twice (with trends).

\vspace{0.45em}

\noindent
\panelbox{3.2cm}{2.4cm}{pooled}\hfill
\panelbox{3.2cm}{2.4cm}{split by group}
}{0.5cm}{4.1cm}

\vfill

% ── 6. Two variables, two types ──────────────────────────────────────────────
\guidedopen{6. Choosing a Plot: Two Variables}{
Name and sketch a plot for data with \textbf{two~categorical} variables.
(And for \textbf{one quantitative and one categorical} variable, if you want)
}{2.6cm}{4.2cm}

\vfill

% ── 7. Proportional ink ──────────────────────────────────────────────────────
\guidedopen{7. The Principle of Proportional Ink}{
Write down the principle.
}{2.8cm}{4.3cm}

\vfill

\end{document}
